\section{Modelo conceptual institucional}
El modelo de información integra dominios documentales académicos y administrativos, garantizando consistencia semántica y trazabilidad de metadatos.

\section{Entidades nucleares}
\begin{itemize}
  \item Usuario: identidad de acceso y perfil operativo.
  \item Archivo Institucional: unidad documental académica con metadatos descriptivos.
  \item Persona RRHH: entidad titular del expediente laboral.
  \item Archivo RRHH: documento asociado al expediente de una persona.
  \item Taxonomía/Tesauro: términos controlados para clasificación y recuperación.
  \item Evento de auditoría: huella de acciones críticas y evidencia de operación.
\end{itemize}

\section{Gobierno de metadatos}
\begin{itemize}
  \item Normalización de campos obligatorios y opcionales.
  \item Catalogación semántica por tipología y tesauro.
  \item Reglas de calidad para integridad, completitud y unicidad.
  \item Políticas de versionado de estructuras de datos.
\end{itemize}

\section{Reglas de consistencia}
\begin{enumerate}
  \item Todo documento debe mantener contexto de origen, tipo y ubicación.
  \item Los expedientes RRHH deben ser trazables a persona titular.
  \item La clasificación documental debe alinearse con taxonomías vigentes.
  \item Los eventos críticos deben registrar actor, acción, módulo y marca temporal.
\end{enumerate}

\section{Estrategia de evolucion de datos}
\begin{itemize}
  \item Migración progresiva a esquema relacional empresarial.
  \item Endurecimiento de validaciones de integridad referencial.
  \item Preparación para búsqueda semántica y analítica histórica.
\end{itemize}
